\documentclass[11pt,a4paper]{article}

%% PACHETE MATEMATICE
\usepackage{amsmath,amssymb,amsfonts}
\usepackage{amsthm}
\usepackage{mathpazo}
\usepackage{eulervm}
\usepackage{enumerate}
\usepackage{color}
\usepackage{graphicx}
\usepackage[all]{xy}
\usepackage{xcolor}
\usepackage{framed}
\usepackage[unicode]{hyperref}
\setcounter{MaxMatrixCols}{20}
\usepackage{multicol}
% \usepackage[romanian]{babel}


%% ASPECT
\linespread{1.05}
\usepackage[T1]{fontenc}
\usepackage[utf8x]{inputenc}
\usepackage[margin=2cm]{geometry}
% \usepackage[color]{showkeys}
\usepackage{anyfontsize}
\usepackage{titlesec}
\titleformat*{\section}{\large\bfseries}

\usepackage{fancyhdr}
\pagestyle{fancy}
\rhead{\small \color{gray} Notițe de Adrian Manea}
\lhead{\small \color{gray} ETTI-AM2, 2017---2018, M. Joița \& A. Niță}


\usepackage[autostyle = false, style = german]{csquotes}
\usepackage[romanian]{babel}
\newcommand\qq{\enquote}

%% COMENZILE MELE
\renewcommand*{\proofname}{Dem.:}
\newcommand\bb{\mathbb}
\newcommand\dr{\mathrm}
\newcommand\kal{\mathcal}
\newcommand\fk{\mathfrak}
\newcommand\op{\oplus}
\newcommand\ot{\otimes}
\newcommand\ds{\displaystyle}
\newcommand\id{\indent}
\newcommand\nid{\noindent}
\newcommand\seq{\subseteq}
\newcommand\sneq{\subsetneq}
\newcommand\speq{\supseteq}
\newcommand\spneq{\supsetneq}
\newcommand\rar{\rightarrow}
\newcommand\Rar{\Rightarrow}
\newcommand\xrar{\xrightarrow}
\newcommand\lar{\leftarrow}
\newcommand\Lar{\Leftarrow}
\newcommand\xlar{\xleftarrow}
\newcommand\lrar{\leftrightarrow}
\newcommand\Lrar{\Leftrightarrow}
\newcommand\ideal{\trianglelefteq}
\newcommand\ai{\text{ a.\^{i}. }}
\newcommand\sk{(\textit{Schi\c{t}\u{a}}) }
\newcommand\rhup{\rightharpoonup}
\newcommand\rhdn{\rightharpoondown}
\newcommand\lhup{\leftharpoonup}
\newcommand\lhdn{\leftharpoondown}
\newcommand\vid{\emptyset}
\newcommand\wh{\widehat}
\newcommand\ol{\overline}
\newcommand\NN{\mathbb{N}}
\newcommand\RR{\mathbb{R}}
\newcommand\ZZ{\mathbb{Z}}
\newcommand\QQ{\mathbb{Q}}
\newcommand\CC{\mathbb{C}}

%% STILURILE MELE
\newtheoremstyle{dotless}{}{}{\itshape}{}{\bfseries}{:}{ }{}

\theoremstyle{dotless}
\newtheorem{theorem}{Teorem\u{a}}[section]
\newtheorem{proposition}{Propozi\c{t}ie}[section]
\newtheorem{lemma}{Lem\u{a}}[section]
\newtheorem{corollary}{Corolar}[section]
\newtheorem{exercise}{Exerci\c{t}iu}

\newtheoremstyle{dotlessdr}{}{}{}{}{\bfseries}{:}{ }{}
\theoremstyle{dotlessdr}
\newtheorem{example}{Exemplu}[section]
\newtheorem{definition}{Defini\c{t}ie}[section]
\newtheorem{remark}{Observa\c{t}ie}[section]




\begin{document}
% \definecolor{labelkey}{RGB}{222,0,0}

\begin{center}
  {\large\textbf{Seminar 6 \\
      Ecuații cu derivate parțiale \\
      de ordinul al doilea}}
\end{center}

\vspace{2cm}

\section{Clasificarea și aducerea la forma canonică}

\begin{definition}\label{def:edp2}
  Se numește \emph{ecuație cvasiliniară cu derivate parțiale de ordinul II},
  cu două variabile independente, o ecuație de forma:
  \begin{equation}
    \label{eq:edp2}
    A(x, y) \frac{\partial^2 z}{\partial x^2} + %
    2B(x, y) \frac{\partial^2 z}{\partial x \partial y} + %
    C(x, y) \frac{\partial^2 z}{\partial y^2} + %
    D\Big(x, y, z, \frac{\partial z}{\partial x}, \frac{\partial z}{\partial y}\Big) = 0,
  \end{equation}
  unde $ A, B, C $ sînt funcții reale, continue pe un deschis din $ \RR^2 $, iar funcția
  $ D $ este continuă.
\end{definition}

\begin{definition}\label{def:curbe-car}
  Se numesc \emph{curbe caracteristice} pentru ecuația (\ref{eq:edp2}) curbele care
  se află pe suprafețele integrale ale ecuației, ale căror proiecții pe planul $ XOY $
  satisfac \emph{ecuația caracteristică}:
  \begin{equation}
    \label{eq:ec-car}
    A(x, y) dy^2 - 2B(x, y) dxdy + C(x, y) dx^2 = 0.
  \end{equation}
\end{definition}

Clasificarea ecuațiilor se face în funcție de curbele caracteristice. Astfel, avem:
\begin{itemize}
\item Dacă $ AC - B^2 < 0 $, ecuația este \emph{de tip hiperbolic};
\item Dacă $ AC - B^2 = 0 $, ecuația este \emph{de tip parabolic};
\item Dacă $ AC - B^2 > 0 $, ecuația este \emph{de tip eliptic}.
\end{itemize}

%%%%%%%%%%%%%%%%%%%%%%%%%%%%%%%%%%%%%%%%%%%%%%%%%%%%%%%%%%%%%%%%%%%%%%

\section{Exemple din fizica matematică}

Unele dintre cele mai importante exemple de ecuații de ordinul al doilea cu
derivate parțiale provin din fizica matematică.

\textbf{Ecuația coardei vibrante}, care are aceeași formă cu \textbf{ecuația undelor plane}:
\[
  \frac{\partial^2 u}{\partial x^2} - \frac{1}{a^2} \frac{\partial^2 u}{\partial t^2} = 0, %
  \quad a^2 = \frac{\rho}{T_0},
\]
unde $ \rho $ este densitatea liniară a coardei, iar $ T_0 $ este tensiunea la care este
supusă coarda în poziția de repaus.

Aceasta este o ecuație de tip hiperbolic.

\vspace{1cm}

\textbf{Ecuația căldurii}, care are forma:
\[
  \frac{\partial^2 u}{\partial x^2} = \frac{1}{a^2} \frac{\partial u}{\partial t}, %
  a^2 = \frac{k}{c \rho},
\]
unde $ k $ este coeficientul de conductibilitate termică, $ c $ este căldura
specifică, iar $ \rho $ este densitatea.

Această ecuație este de tip parabolic.

\vspace{1cm}

\textbf{Ecuația lui Laplace}, cu forma generală:
\[
  \frac{\partial^2 u}{\partial x^2} + \frac{\partial^2 u}{\partial y^2} = 0.
\]

Să remarcăm că aceasta este echivalentă cu $ \Delta u = 0 $, adică $ u $ să fie
funcție armonică (lucru care justifică și numele ecuației).

Ecuația lui Laplace este de tip eliptic.

%%%%%%%%%%%%%%%%%%%%%%%%%%%%%%%%%%%%%%%%%%%%%%%%%%%%%%%%%%%%%%%%%%%%%%

\section{Direcții caracteristice și forma canonică}

\begin{definition}\label{def:dir-car}
  Ecuațiile de forma:
  \begin{equation}
    \label{eq:dir-car}
    \frac{dy}{dx} = \mu_1(x, y), \quad \frac{dy}{dx} = \mu_2(x, y)
  \end{equation}
  asociate ecuației curbelor caracteristice (\ref{eq:ec-car}) se numesc
  \emph{curbe caracteristice} ale ecuației (\ref{eq:edp2}).
\end{definition}

Prin integrarea celor două ecuații (\ref{eq:dir-car}) se obțin două familii de
curbe în planul $ XOY $, de forma $ \varphi_1(x, y) = c_1 $ și $ \varphi_2(x, y) = c_2 $,
cu $ c_1, c_2 $ constante arbitrare, curbe care sînt proiecțiile pe planul $ XOY $ ale
curbelor caracteristice.

Pentru a aduce ecuațiile (\ref{eq:edp2}) în forma canonică, distingem
următoarele cazuri:
\begin{itemize}
\item Dacă ecuația este de \emph{tip hiperbolic}, facem schimbarea de variabile
  $ \tau = \varphi_1(x, y) $ și $ \eta = \varphi_2(x, y) $ și \emph{prima forma canonică} se
  obține a fi:
  \[
    \frac{\partial^2 z}{\partial \tau \partial \eta} + %
    \Psi_1\Big( \tau, \eta, z, \frac{\partial z}{\partial \tau}, \frac{\partial z}{\partial \eta} \Big) = 0.
  \]

  Cu transformarea $ \tau = x + y $ și $ \eta = x - y $ se obține, din ecuația anterioară,
  \emph{a doua formă canonică}, anume:
  \[
    \frac{\partial^2 z}{\partial x^2} - \frac{\partial^2 z}{\partial y^2} + %
    \Psi'_1 \Big( \tau, \eta, z, \frac{\partial z}{\partial x}, \frac{\partial z}{\partial y} \Big) = 0.
  \]
\item Dacă ecuația este de \emph{tip parabolic}, avem
  $ \varphi_1(x, y) = \varphi_2(x, y) = \varphi(x, y) $ și putem face schimbarea de
  variabile $ \tau = \varphi(x, y) $ și $ \eta = x $, ajungînd la forma canonică:
  \[
    \frac{\partial^2 z}{\partial \eta^2} + %
    \Psi_2\Big( \tau, \eta, z, \frac{\partial z}{\partial \tau}, \frac{\partial z}{\partial \eta} \Big) = 0.
  \]
\item Pentru ecuații de \emph{tip eliptic}, funcțiile $ \varphi_1(x, y) $ și $ \varphi_2(x, y) $
  sînt complex conjugate și putem nota $ \alpha(x, y) = \mathrm{Re} \varphi_1(x, y) $, iar
  $ \beta(x, y) = \mathrm{Im} \varphi_1(x, y) $. Atunci, cu schimbarea de variabile de forma
  $ \tau = \alpha(x, y) $ și $ \eta = \beta(x, y) $, ajungem la forma canonică:
  \[
    \frac{\partial^2 z}{\partial \tau^2} + \frac{\partial^2 z}{\partial \eta^2} + %
    \Psi_3 \Big( \tau, \eta, z, \frac{\partial z}{\partial \tau}, \frac{\partial z}{\partial \eta} \Big) = 0.
  \]
\end{itemize}

%%%%%%%%%%%%%%%%%%%%%%%%%%%%%%%%%%%%%%%%%%%%%%%%%%%%%%%%%%%%%%%%%%%%%%

\section{Cazul coeficienților constanți}

Dacă avem o ecuație liniară și omogenă în raport cu derivatele parțiale de ordinul
al doilea, cu coeficienți constanți:
\[
  A \frac{\partial^2 z}{\partial x^2} + 2B \frac{\partial^2 z}{\partial x \partial y} + %
  C \frac{\partial^2 z}{\partial y^2} = 0,
\]
cu $ A, B, C $ constante, atunci ecuația diferențială a proiecțiilor curbelor
caracteristice pe planul $ XOY $ are forma particulară:
\[
  A dy^2 - 2B dxdy + C dx^2 = 0.
\]

Direcțiile caracteristice sînt date, atunci, de:
\[
  \begin{cases}
    dy - \mu_1 dx &= 0 \\
    dy - \mu_2 dx &= 0
  \end{cases}
\]
care dau simplu soluția:
\[
  \begin{cases}
    y - \mu_1 x &= c_1 \\
    y - \mu_2 x &= c_2
  \end{cases}, c_1, c_2 \in \RR.
\]

Aducerea la forma canonică se face mai simplu:
\begin{itemize}
\item Dacă ecuația este \emph{de tip hiperbolic}, facem schimbarea de variabile:
  \[
    \begin{cases}
      \tau &= y - \mu_1 x \\
      \eta &= y - \mu_2 x
    \end{cases},
  \]
  iar ecuația devine:
  \[
    \frac{\partial^2 z}{\partial \tau \partial \eta} = 0,
  \]
  care are soluția generală $ z = f(\tau) + g(\eta) $, unde $ f, g $ sînt funcții arbitrare.
  Înlocuind pentru a reveni la vechile variabile, avem:
  \[
    z(x, y) = f(y - \mu_1 x) + g(y - \mu_2 x).
  \]
\item Pentru \emph{cazul parabolic}, avem $ \mu_1 = \mu_2 = \frac{B}{A} $, iar ecuația
  curbelor devine atunci $ A dy - B dx = 0 $, cu integrala generală $ Ay - Bx = c \in \RR $.
  Schimbarea de variabile $ \tau = Ax - By$ și $ \eta = x $ conduce la forma canonică:
  \[
    \frac{\partial^2 z}{\partial \eta^2} = 0.
  \]
  Soluția generală este $ z = \eta f(\tau) + g(\tau) $, unde $ f, g $ sînt funcții arbitrare.
\item Pentru \emph{cazul eliptic}, forma canonică este chiar ecuația Laplace:
  \[
    \frac{\partial^2 z}{\partial \tau^2} + \frac{\partial^2 z}{\partial \eta^2} = 0.
  \]
\end{itemize}

%%%%%%%%%%%%%%%%%%%%%%%%%%%%%%%%%%%%%%%%%%%%%%%%%%%%%%%%%%%%%%%%%%%%%%

\section{Exerciții}

1. Aduceți la forma canonică următoarele ecuații:
\begin{enumerate}[(a)]
\item $ \frac{\partial^2 u}{\partial x^2} + 2 \frac{\partial^2 u}{\partial x \partial y} - %
  3 \frac{\partial^2 u}{\partial y^2} + 2 \frac{\partial u}{\partial x} + %
  6 \frac{\partial u}{\partial y} = 0 $;
\item $ 3 \frac{\partial^2 u}{\partial x^2} + 7 \frac{\partial^2 u}{\partial x\partial y} + %
  2 \frac{\partial^2 u}{\partial y^2} = 0 $;
\item $ 4 \frac{\partial^2 u}{\partial x^2} + 4 \frac{\partial^2 u}{\partial x \partial y} + %
  \frac{\partial^2 u}{\partial y^2} - 2 \frac{\partial u}{\partial y} = 0 $;
\item $ \frac{\partial^2 u}{\partial x^2} - 6 \frac{\partial^2 u}{\partial x\partial y} + %
  10 \frac{\partial^2 u}{\partial y^2} + \frac{\partial u}{\partial x} - %
  3 \frac{\partial u}{\partial y} = 0 $;
\item $ 2 \frac{\partial^2 u}{\partial x^2} - 7 \frac{\partial^2 u}{\partial x \partial y} + %
  3 \frac{\partial^2 u}{\partial y^2} = 0 $;
\item $ y \frac{\partial^2 u}{\partial x^2} + (x + y) \frac{\partial^2 u}{\partial x \partial y} + %
  x \frac{\partial^2 u}{\partial y^2} = 0 $;
\item $ (1 + x^2) \frac{\partial^2 u}{\partial x^2} + %
  (1 + y^2) \frac{\partial^2 u}{\partial y^2} + %
  x \frac{\partial u}{\partial x} + y \frac{\partial u}{\partial y} = 0 $;
\item $ x^2 \frac{\partial^2 u}{\partial x^2} - %
  2xy \cdot \frac{\partial^2 u}{\partial x \partial y} + %
  y^2 \cdot \frac{\partial^2 u}{\partial y^2} + %
  x \cdot \frac{\partial u}{\partial x} + y \cdot \frac{\partial u}{\partial y} = 0 $;
\end{enumerate}

\emph{Soluție: (a)} Deoarece avem $ A = 1, B = 1, C = -3 $, rezultă $ B^2 - AC = 4 > 0 $,
deci ecuația este de tip hiperbolic.

Scriem ecuația caracteristică:
\[
  \Big( \frac{dy}{dx} \Big)^2 - 2 \frac{dy}{dx} - 3 = 0,
\]
care poate fi rezolvată ca o ecuație de gradul al doilea, de unde:
\[
  \begin{cases}
    \frac{dy}{dx} = 3 &\Rightarrow y - 3x = c_1 \\
    \frac{dy}{dx} = -1 &\Rightarrow y + x = c_2
  \end{cases}.
\]

Cu schimbarea de variabile:
\[
  \begin{cases}
    \tau &= y - 3x \\
    \eta &= y + x
  \end{cases},
\]
funcția căutată $ u(x, y) $ devine $ u(\tau, \eta) $, astfel că toate derivatele parțiale
se calculează acum folosind formula funcțiilor implicite și derivarea funcțiilor compuse. Obținem, succesiv:
\begin{align*}
  \frac{\partial u}{\partial x} &= \frac{\partial u}{\partial \tau} \cdot \frac{\partial \tau}{\partial x} + %
                                  \frac{\partial u}{\partial \eta} \cdot \frac{\partial \eta}{\partial x} \\
                                &=- 3\frac{\partial u}{\partial \tau} + \frac{\partial u}{\partial \eta} \\
  \frac{\partial u}{\partial y} &= \frac{\partial u}{\partial \tau} \cdot \frac{\partial \tau}{\partial y} + %
                                  \frac{\partial u}{\partial \eta} \cdot \frac{\partial \eta}{\partial y} \\
                                &=\frac{\partial u}{\partial \tau} + \frac{\partial u}{\partial \eta} \\
  \frac{\partial^2 u}{\partial x^2} &= \frac{\partial}{\partial x} \Big( \frac{\partial u}{\partial x} \Big)  \\
                                &= \frac{\partial}{\partial x} \Big( \frac{\partial u}{\partial \tau} \cdot %
                                  \frac{\partial \tau}{\partial x} + \frac{\partial u}{\partial \eta} \cdot \frac{\partial \eta}{\partial x} \Big) \\
                                &= \frac{\partial}{\partial x} \Big( \frac{\partial u}{\partial \tau} \cdot \frac{\partial \tau}{\partial x}\Big) + %
                                  \frac{\partial}{\partial x} \Big( \frac{\partial u}{\partial \eta} \cdot \frac{\partial \eta}{\partial x} \Big) \\
                                &= \frac{\partial}{\partial x} \Big( \frac{\partial u}{\partial \tau} \Big) \cdot %
                                  \frac{\partial \tau}{\partial x} + \frac{\partial u}{\partial \tau} \cdot %
                                  \frac{\partial^2 \tau}{\partial x^2} + \frac{\partial}{\partial x} \Big( \frac{\partial u}{\partial \eta} \Big) \cdot %
                                  \frac{\partial \eta}{\partial x} + \frac{\partial u}{\partial \eta} \cdot 5
                                  \frac{\partial^2 \eta}{\partial x^2} \\
                                &= \Big[ \frac{\partial}{\partial \tau} \Big(\frac{\partial u}{\partial \tau}\Big) \cdot %
                                  \frac{\partial \tau}{\partial x} + \frac{\partial}{\partial \eta} \Big(\frac{\partial u}{\partial \tau}\Big) \cdot %
                                  \frac{\partial \eta}{\partial x} \Big] \cdot \frac{\partial \tau}{\partial x} + \\
                                &+ \Big[ \frac{\partial}{\partial \eta} \Big(\frac{\partial u}{\partial \eta}\Big) \cdot %
                                  \frac{\partial \eta}{\partial x} + \frac{\partial}{\partial \tau} \Big(\frac{\partial u}{\partial \eta}\Big) \cdot %
                                  \frac{\partial \eta}{\partial x} \Big] \cdot \frac{\partial \eta}{\partial x} + \\
                                &+ \frac{\partial^2 \tau}{\partial x^2} \cdot \frac{\partial u}{\partial \tau} + %
                                  \frac{\partial u}{\partial \eta} \cdot \frac{\partial^2 \eta}{\partial x^2} \\
                                &= \frac{\partial^2 u}{\partial \tau^2} \cdot \Big( \frac{\partial \tau}{\partial x} \Big)^2 + %
                                  \frac{\partial u}{\partial \eta \partial \tau} \cdot \frac{\partial \eta}{\partial x} \cdot \frac{\partial \tau}{\partial x} + %
                                  \frac{\partial^2 u}{\partial^2 \eta} \cdot \Big( \frac{\partial \eta}{\partial x} \Big)^2 + \\
                                &+ \frac{\partial^2 u}{\partial \tau\partial \eta} \cdot \frac{\partial \tau}{\partial x} \cdot \frac{\partial \eta}{\partial x} + %
                                  \frac{\partial^2 \tau}{\partial x^2} \cdot \frac{\partial u}{\partial \tau} + \frac{\partial u}{\partial \eta} \cdot %
                                  \frac{\partial^2 \eta}{\partial x^2} \\
                                &= 9 \frac{\partial^2 u}{\partial \tau^2} - %
                                      6 \frac{\partial^2 u}{\partial \tau \partial \eta} + %
                                      \frac{\partial^2 u}{\partial \eta^2} \\
  \frac{\partial^2 u}{\partial x \partial y} &= \frac{\partial}{\partial x}\Big( \frac{\partial u}{\partial y} \Big)\\
                                &= -3\frac{\partial^2 u}{\partial \tau^2} - %
                                               2 \frac{\partial ^2 u}{\partial \tau \partial \eta} + %
                                               \frac{\partial^2 u}{\partial \eta^2} \\
  \frac{\partial^2 u}{\partial y^2} &= \frac{\partial^2 u}{\partial \tau^2} + %
                                      2 \frac{\partial^2 u}{\partial \tau \partial \eta} + %
                                      \frac{\partial^2 u}{\partial \eta^2}.  
\end{align*}

Toate derivatele parțiale în raport cu $ x $ și $ y $ se calculează, deci, ținînd cont de legătura
cu noile variabile $ \eta $ și $ \tau $. Practic, avem:
\[
  \frac{\partial}{\partial x} = \frac{\partial}{\partial \eta} \cdot \frac{\partial \eta}{\partial x} + %
  \frac{\partial}{\partial \tau} \cdot \frac{\partial \tau}{\partial x}
\]
și similar pentru $ y $.

Rezultă că, în final, forma canonică este:
\[
  -16 \frac{\partial^2 u}{\partial \tau \partial \eta} + 8 \frac{\partial u}{\partial \eta} = 0 %
  \Leftrightarrow \frac{\partial^2 u}{\partial \tau \partial \eta} - %
  \frac{1}{2} \frac{\partial u}{\partial \eta} = 0.
\]

\vspace{1cm}

2. Determinați soluția ecuației:
\[
  \frac{\partial^2 u}{\partial x^2} + 2 \frac{\partial^2 u}{\partial x \partial y} - %
  3 \frac{\partial^2 u}{\partial y^2} = 0,
\]
care satisface condițiile:
\[
  \begin{cases}
    u(x, y) &= 3x^2 \\
    \frac{\partial u}{\partial y}(x, 0) &= \cos x
  \end{cases}.
\]

\vspace{1cm}

3. Rezolvați ecuația:
\[
  3 \frac{\partial^2 u}{\partial x^2} + 7 \frac{\partial^2 u}{\partial x\partial y} + %
  2 \frac{\partial^2 u}{\partial y^2} = 0,
\]
cu condițiile:
\[
  \begin{cases}
    u(x, y) &= x^3 \\
    \frac{\partial u}{\partial y}(x, 0) &= 2x^2
  \end{cases}.
\]

\end{document}